%! Introducing Vectors — Unit 1, Lesson 1.0 — Slides (Beamer)
\documentclass[aspectratio=169, 11pt]{beamer}
\usepackage{linalg-beamer}   % provides \burgundyheader{} and \sectionlabel[color]{}
\newcommand{\vv}{\mathbf{v}}
\newcommand{\ww}{\mathbf{w}}

\begin{document}

% TITLE SLIDE (CourseName is not defined in beamer, so the name is literal)
{
\setbeamercolor{background canvas}{bg=burgundy}
\begin{frame}[plain]
  \vspace{2.8cm}\hspace{0.55cm}%
  \begin{minipage}{0.9\textwidth}
    {\color{goldacc}\bfseries\small LESSON 1.0}\\[0.5cm]
    {\color{white}\bfseries\LARGE Introducing Vectors}\\[0.3cm]
    {\color{white}\bfseries\large Vocabulary and Length}\\[1.0cm]
    {\color{blushmid}\small Linear Algebra~$\cdot$~Unit 1: Vectors and Matrices}
  \end{minipage}
\end{frame}
}

% HOOK
\begin{frame}[t, plain]
  \burgundyheader{A trip in two numbers}
  \sectionlabel{Hook}
  \begin{columns}[T]
    \begin{column}{0.58\textwidth}
      A drone flies \textbf{3 blocks east} and \textbf{4 blocks north}.\\[6pt]
      The single pair $(3,4)$ records the \emph{whole} trip.\\[10pt]
      \textbf{Question.} How far is the drone from where it started, in a straight line?
    \end{column}
    \begin{column}{0.40\textwidth}
      \centering
      \begin{tikzpicture}[scale=0.5]
        \draw[step=1, black!15, thin] (-0.3,-0.3) grid (4.3,4.3);
        \draw[->, thick, black!60] (0,0)--(4.4,0);
        \draw[->, thick, black!60] (0,0)--(0,4.4);
        \draw[->, very thick, burgundy] (0,0)--(3,4) node[midway, above left]{$(3,4)$};
      \end{tikzpicture}
    \end{column}
  \end{columns}
\end{frame}

% WHAT IS A VECTOR
\begin{frame}[t, plain]
  \burgundyheader{What is a vector?}
  \sectionlabel{Definition}
  A \textbf{vector} is an ordered list of numbers. In the plane it has two
  \textbf{components}:
  \[
    \vv = \begin{bmatrix} v_1 \\ v_2 \end{bmatrix}
        = \begin{bmatrix} 3 \\ 4 \end{bmatrix}.
  \]
  \begin{itemize}
    \item Picture it as a \textbf{point} $(v_1,v_2)$, or as an \textbf{arrow} from the origin.
    \item First component $=$ how far \textbf{east}; second $=$ how far \textbf{north}.
    \item A single number (like $5$) is a \textbf{scalar}: size, but no direction.
  \end{itemize}
\end{frame}

% ADDING AND SCALING
\begin{frame}[t, plain]
  \burgundyheader{Adding and scaling}
  \sectionlabel{Operations}
  \begin{columns}[T]
    \begin{column}{0.5\textwidth}
      \textbf{Add} component by component:
      \[
        \vv+\ww=\begin{bmatrix} v_1+w_1 \\ v_2+w_2 \end{bmatrix}
      \]
      Geometrically: arrows \textbf{tip to tail}.
    \end{column}
    \begin{column}{0.5\textwidth}
      \textbf{Scale} every component by $c$:
      \[
        c\,\vv=\begin{bmatrix} c\,v_1 \\ c\,v_2 \end{bmatrix}
      \]
      $c>1$ longer, $0<c<1$ shorter, $c<0$ flips.
    \end{column}
  \end{columns}
  \vspace{10pt}
  \begin{block}{The big idea of Unit 1}
    Scaling \emph{and} adding together, $c\,\vv + d\,\ww$, is a \textbf{linear combination}.
  \end{block}
\end{frame}

% LENGTH
\begin{frame}[t, plain]
  \burgundyheader{How long is a vector?}
  \sectionlabel[goldacc]{Magnitude}
  The \textbf{length} (magnitude) comes straight from the Pythagorean theorem:
  \[
    \|\vv\| = \sqrt{\,v_1^{\,2}+v_2^{\,2}\,}.
  \]
  \begin{itemize}
    \item Example: $\left\|\begin{bsmallmatrix} 3\\4 \end{bsmallmatrix}\right\|
          = \sqrt{9+16} = \sqrt{25} = 5$ --- the drone is $5$ blocks away.
    \item Length is a \textbf{scalar} and is \textbf{never negative}.
    \item A \textbf{unit vector} has length $1$: shrink $\vv$ by dividing by $\|\vv\|$.
  \end{itemize}
\end{frame}

% WRAP / PREVIEW
\begin{frame}[t, plain]
  \burgundyheader{Wrap-up and what's next}
  \sectionlabel{Connections}
  \textbf{Today:} vectors as components and arrows; adding, scaling, and length.\\[8pt]
  \textbf{Next --- Lesson 1.1:} do both operations at once, $c\,\vv + d\,\ww$, and ask
  \emph{which points in the plane you can reach} by combining two given vectors.
\end{frame}

\end{document}
